% This Source Code Form is subject to the terms of the Mozilla Public
% License, v. 2.0. If a copy of the MPL was not distributed with this
% file, You can obtain one at https://mozilla.org/MPL/2.0/.

\documentclass[a4paper,12pt]{article}
\usepackage[utf8]{inputenc}
\usepackage[russian]{babel}
\usepackage{amsmath, amssymb}
\usepackage[hyperfootnotes=false]{hyperref}
\usepackage{geometry}
\usepackage{indentfirst}
\usepackage{graphicx}
\usepackage{float}
\usepackage{enumitem}
\usepackage{caption}
\usepackage{array}
\geometry{
	a4paper,
	top=24mm,
	bottom=30mm,
	left=20mm,
	right=20mm
}

\newcommand{\blfootnote}[1]{%
	\begingroup
	\renewcommand{\thefootnote}{}%
	\footnote{#1}%
	\addtocounter{footnote}{-1}%
	\endgroup
}

% Отступ первой строки абзаца 1 см
\setlength{\parindent}{1cm}
\setlength{\parskip}{0pt}
\linespread{1}                 % одинарный интервал

% Отключение нумерации страниц
\pagenumbering{gobble}

% Настройка подписей рисунков и таблиц
% Используем \caption* для отключения точки, но оставляем обычный \caption
\usepackage[labelsep=space, justification=centering]{caption}
\captionsetup[figure]{font=small}   % small = 11pt при 12pt основного шрифта
\captionsetup[table]{font=small}
\renewcommand{\figurename}{Рис.}
\renewcommand{\tablename}{Табл.}

% Убираем точку в конце подписи (переопределяем \caption)
\usepackage{etoolbox}
\makeatletter
\renewcommand{\fnum@figure}{\figurename~\thefigure}
\renewcommand{\fnum@table}{\tablename~\thetable}
\apptocmd{\@caption}{\def\@captionheadfont{\normalfont}\def\@captionfont{\normalfont}}{}{}
\makeatother

% Настройка списка литературы
\makeatletter
\renewcommand{\@biblabel}[1]{#1.}   % нумерация "1.", "2."
\makeatother

\begin{document}
	\subsubsection*{Classical trajectory mеthod. Quantum Entanglement from Classical Trajectories}
	
	Today I will talk about the classical trajectories method. 
	We use this method to simulate large molecules or solids. 
	The idea is simple: nuclei move by Newton's laws (like small balls), 
	and electrons are treated quantum mechanically.
	
	\textbf{Example: sodium iodide molecule (NaI).}
	
	In the ground state, the molecule is Na\(^+\) and I\(^-\). 
	Nuclei move classically. 
	The laser gives energy to the electron, and it jumps to another state. 
	Now the molecule is Na I. It starts to fly apart.
	
	Near the point where the two energy levels cross, 
	the electron can jump back or stay. 
	Two outcomes are possible:
	\begin{itemize}
		\item ions: Na\(^+\) + I\(^-\)
		\item neutral atoms: Na + I
	\end{itemize}
	
	A good classical trajectories method must show this splitting (branching).
	
	\textbf{General approach (mixed quantum-classical dynamics).}
	
	In standard mixed methods, electrons are described quantum mechanically — 
	by the equation for the density matrix (Liouville–von Neumann equation). 
	This equation is solved together with Newton's equations for the nuclei.
	
	\begin{equation}
		i\hbar \frac{\partial \rho}{\partial t} = \left[\hat{H}, \rho\right]
	\end{equation}
	
	We can add extra terms:
	\begin{itemize}
		\item relaxation term — describes loss of coherence,
		\item pump term — describes the laser action (through the dipole moment).
	\end{itemize}
	
	\textbf{What the authors did.}
	
	Runeson and Richardson showed that quantum entanglement can be described 
	using ordinary classical trajectories. No jumps, no randomness, no coupled trajectories. 
	All quantum effects are hidden in the initial weights of the trajectories. 
	These weights can be negative — and that is what creates quantum interference.
	
	\textbf{Mapping a two-level system to a classical spin.}
	
	A two-level system (like an electron) is like a spin 1/2. 
	The authors map it onto a classical vector \(s\). 
	But the length of this vector is not \(1/2\). It is \(\sqrt{3}/2\). 
	This is not the usual Bloch sphere. It is called the Stratonovich–Weyl representation.
	
	\textbf{Path of N spins and the centroid.}
	
	Instead of one spin, the authors take \(N\) classical spins 
	\(s_1, s_2, \dots, s_N\). 
	Then they compute their average — the centroid \(\bar{s} = \frac{1}{N}\sum_k s_k\). 
	Surprisingly, quantum observables depend only on this centroid, 
	not on the individual spins.
	
	\textbf{Universal weight function \(G_N(\bar{s})\).}
	
	After integrating over all spins, only the centroid remains. 
	The result is a universal weight function \(G_N(\bar{s})\). 
	It depends only on the length \(|\bar{s}|\). 
	For \(N=1\), the centroid is always on a sphere of radius \(\sqrt{3}/2\). 
	For \(N=4,8\), the distribution of \(\bar{s}\) concentrates near \(\pm 1/2\) — 
	the quantum eigenvalues. 
	There are also negative regions in \(G_N\). 
	These negative weights create quantum interference.
	
	\textbf{Dynamics is fully classical.}
	
	The centroid \(\bar{s}\) precesses like a classical spin: 
	\(\dot{\bar{s}} = \bar{s} \times \mathbf{H}(x)\). 
	The nuclear coordinates \(x\) and momenta \(p\) follow Newton's laws. 
	Everything is deterministic. 
	No stochastic jumps. 
	The weights \(G_N\) are conserved along each trajectory.
	
	\textbf{Tully model I — wave packet branching.}
	
	A wave packet enters an avoided crossing. 
	It must split into two parts. 
	Ehrenfest dynamics gives only one peak (average path). 
	With \(N=1\) (linearized spin mapping), we get an envelope — but still one peak. 
	With \(N=4\) or \(N=8\), two clear peaks appear. 
	The method correctly reproduces the quantum branching.
	
	\textbf{Tully model III — recrossing.}
	
	The wave packet passes through the coupling region twice. 
	Surface hopping (FSSH) gives wrong oscillations (overcoherence). 
	The spin path-integral method with \(N=4\) gives the correct result — 
	no oscillations, correct entanglement entropy.
	
	\textbf{Main conclusion.}
	
	Quantum entanglement can be described by an ensemble of classical trajectories. 
	To do this, we need:
	\begin{itemize}
		\item many classical spins,
		\item their average (centroid),
		\item initial weights \(G_N\) that can be negative.
	\end{itemize}
	The dynamics is classical. 
	The method works where other methods fail.
	
\end{document}